\documentclass[11pt]{article}

\usepackage[margin=1in]{geometry}
\usepackage[T1]{fontenc}
\usepackage[utf8]{inputenc}
\usepackage{lmodern}
\usepackage{microtype}
\usepackage{amsmath,amssymb}
\usepackage{booktabs}
\usepackage{xcolor}
\usepackage{hyperref}
\usepackage{enumitem}
\usepackage{graphicx}
\usepackage{float}

\hypersetup{colorlinks=true,linkcolor=blue!50!black,urlcolor=blue!50!black,citecolor=blue!50!black,hypertexnames=false}

\title{Causal Interaction Memory Toy:\\Phase-Aware Retry and Transfer Without Causal-Model Claims}
\author{VLA Ideas}
\date{September 1, 2026}

\begin{document}
\maketitle

\begin{abstract}
This note distills one idea seeded by ZeVA into a deterministic manipulation-memory toy. Related tasks share hidden mass, friction, and joint-stiffness factors. Methods receive only visible task context plus prior interaction memory from sibling tasks and retry observations from the current target task. A phase-aware dual-timescale memory achieves the best retry success, strongest robustness, and lowest attempts-to-solve, while flat transition retrieval is competitive and lightweight online fine-tuning helps when no related-task memory exists. The result is a mechanism probe only; it does not identify causal world models or reproduce ZeVA.
\end{abstract}

\section{Motivation}
Long-horizon interaction tasks often fail because the visible goal stays similar while the hidden response of the object changes. A heavy drawer, a sticky slider, and a stiff hinge can each demand a different corrective force even when the scene semantics are unchanged. The question here is narrower than ZeVA: if those hidden properties are shared across related tasks, what memory interface best helps retries and transfer?

\section{Seed Reference}
The experiment is seeded by ZeVA~\cite{zeva}, which motivates interaction-driven memory for long-horizon VLA systems. This package keeps only the memory theme and replaces perception, language, and robot execution with an explicit low-dimensional toy. No claim is made that the methods recover identified causes; they only exploit structured correlations induced by hidden physical factors.

\section{Toy Setup}
Each task has three ordered phases:
\begin{itemize}[leftmargin=1.5em]
  \item \textbf{lift:} mostly sensitive to hidden mass,
  \item \textbf{slide:} mostly sensitive to hidden friction,
  \item \textbf{turn:} mostly sensitive to hidden joint stiffness.
\end{itemize}
A family of related tasks shares a latent vector over those properties, with small task-specific variation. For each target task, the method predicts three phase-specific force corrections and may retry from the start for up to four attempts. If one phase fails, the attempt stops and yields a biased but informative transition estimate for that phase.

The methods never receive the latent mass/friction/stiffness values directly. They only use visible task context, related-task support memory, and retry-time interaction summaries.

\section{Compared Methods}
\begin{itemize}[leftmargin=1.5em]
  \item \textbf{Raw history:} retrieve whole support-task traces with weak phase structure.
  \item \textbf{Transition retrieval:} phase-specific nearest-neighbor retrieval over support transitions.
  \item \textbf{Phase-aware dual memory:} slow family consolidation plus a fast retry cache and cross-phase covariance update.
  \item \textbf{Online fine-tune:} ridge-style parametric adaptation using support tasks and retry observations.
\end{itemize}
The dual-memory method is intentionally described as a structured memory baseline, not as a causal discovery algorithm.

\section{Metrics}
\begin{itemize}[leftmargin=1.5em]
  \item \textbf{First-attempt transfer:} fraction of target tasks solved before any target-specific retry update.
  \item \textbf{Retry success:} fraction solved within four total attempts.
  \item \textbf{Retry gain:} retry success minus first-attempt transfer.
  \item \textbf{Mean attempts:} attempts-to-solve, with unsolved tasks assigned the full budget plus one.
  \item \textbf{Robustness:} the same metrics under stronger hidden-parameter and context shifts.
\end{itemize}

\section{Main Results}
The checked artifacts were generated with:
\begin{verbatim}
PYTHONWARNINGS=error \
python causal_interaction_memory/run_causal_interaction_memory.py \
  --seed 29 --trials 240 --support-tasks 4 --max-attempts 4
\end{verbatim}

\begin{table}[H]
\centering
\small
\begin{tabular}{lrrrrr}
\toprule
Method & Transfer & Retry success & Retry gain & Mean attempts & Robust split \\
\midrule
Raw history & 30.8\% & 53.3\% & 22.5 pts & 3.17 & 38.3\% \\
Transition retrieval & 44.6\% & 83.3\% & 38.8 pts & 2.40 & 63.3\% \\
Phase-aware dual memory & \textbf{48.3\%} & \textbf{95.4\%} & \textbf{47.1 pts} & \textbf{1.85} & \textbf{85.0\%} \\
Online fine-tune & 39.6\% & 69.6\% & 30.0 pts & 2.73 & 43.3\% \\
\bottomrule
\end{tabular}
\caption{Aggregate results over 240 target tasks with four related support tasks per family. Robust split means the hardest severity bucket from the main evaluation.}
\label{tab:main}
\end{table}

The dual-memory method solves 79.2\% of tasks by the second attempt and 92.1\% by the third, compared with 57.5\% and 74.6\% for flat transition retrieval. It also reduces mean phase error from 0.0476 on attempt 1 to 0.0276 on the final attempt.

\begin{figure}[H]
  \centering
  \includegraphics[width=0.98\textwidth]{../outputs/headline_metrics.png}
  \caption{Headline transfer, retry, and attempts-to-solve comparison.}
  \label{fig:headline}
\end{figure}

\begin{figure}[H]
  \centering
  \includegraphics[width=0.72\textwidth]{../outputs/retry_curves.png}
  \caption{Cumulative solve rate as the retry budget increases from one to four attempts.}
  \label{fig:retry}
\end{figure}

\section{Sample Efficiency and Transfer}
Support-memory size is swept over $\{0,2,4,8,12\}$ related tasks. Two points matter most:
\begin{itemize}[leftmargin=1.5em]
  \item With only two support tasks, phase-aware dual memory already reaches 98.8\% final success.
  \item With no support tasks, online fine-tuning is the best fallback at 53.8\% because it can still use retries, while raw history falls to 5.6\%.
\end{itemize}
At twelve support tasks, online fine-tuning reaches the best first-attempt transfer at 54.4\%, matching the best transfer band, but it still trails dual memory on final success (75.0\% versus 98.1\%).

\begin{figure}[H]
  \centering
  \includegraphics[width=0.98\textwidth]{../outputs/sample_efficiency.png}
  \caption{Support-memory sweep for final retry success and first-attempt transfer.}
  \label{fig:sample}
\end{figure}

\section{Robustness}
Robustness increases the task-specific hidden shift and visible-context extrapolation severity. At severity $1.50$, final success is 27.5\% for raw history, 57.5\% for transition retrieval, 80.6\% for dual memory, and 36.3\% for online fine-tuning. The ranking therefore holds under stronger support--target mismatch, though all methods degrade.

\begin{figure}[H]
  \centering
  \includegraphics[width=0.98\textwidth]{../outputs/robustness.png}
  \caption{Final retry success and first-attempt transfer under stronger hidden-property shifts.}
  \label{fig:robust}
\end{figure}

\section{Memory Ablations}
\begin{table}[H]
\centering
\small
\begin{tabular}{lrrr}
\toprule
Setting & Success & Transfer & Solved by attempt 2 \\
\midrule
Full dual memory & \textbf{96.7\%} & \textbf{41.7\%} & \textbf{76.1\%} \\
No fast retry memory & 41.7\% & 41.7\% & 41.7\% \\
No slow consolidation & 34.4\% & 0.6\% & 8.3\% \\
No cross-phase retrieval & 96.1\% & 41.7\% & 76.1\% \\
No phase-aware prior & 88.3\% & 15.6\% & 51.1\% \\
\bottomrule
\end{tabular}
\caption{Ablations on 180 held-out target tasks.}
\label{tab:ablations}
\end{table}

Removing the fast retry cache or slow consolidation causes the largest losses, confirming that both timescales matter. Removing phase-aware retrieval has a smaller but still meaningful effect, while cross-phase covariance helps only modestly by trimming the late-attempt tail from 96.1\% to 96.7\%.

\begin{figure}[H]
  \centering
  \includegraphics[width=0.98\textwidth]{../outputs/ablations.png}
  \caption{Dual-memory ablations for final success, transfer, and two-attempt adaptation.}
  \label{fig:ablations}
\end{figure}

\section{Interpretation}
The local evidence supports three bounded claims. First, flat phase-specific transition retrieval is already much better than replaying raw whole-task history. Second, separating slow family consolidation from fast retry memory produces the best retry success and robustness in this structured family-sharing setting. Third, lightweight online fine-tuning is useful when no support memory exists, but it is less robust than explicit memory once related tasks are available.

What the experiment does \emph{not} show is equally important: it does not prove that the learned summaries identify true causes, that the same ranking would hold in visual-language robotics, or that ZeVA's world-model components are reproduced here.

\section{Limitations and Follow-Ups}
\begin{itemize}[leftmargin=1.5em]
  \item The toy exposes low-dimensional interaction summaries rather than images, language, or robot trajectories.
  \item Online fine-tuning is a ridge-style proxy, not adaptation of a pretrained VLA.
  \item The family structure is favorable to memory transfer because related tasks really do share hidden factors.
  \item Cross-phase updates are linear covariance updates; richer latent-state models could matter more in harder settings.
  \item A natural follow-up is to add distractor families with partially shared or misleading structure and test whether consolidation remains selective instead of overconfident.
\end{itemize}

\begin{thebibliography}{9}
\bibitem{zeva} \emph{Zeva: In-Context Causal Learning for Self-Evolving Vision-Language-Action Models}. arXiv:2608.30880, 2026. \url{https://arxiv.org/abs/2608.30880}
\end{thebibliography}

\end{document}
